\documentclass[a4paper, 12pt, twoside, openright]{book}
\usepackage[english,italian]{babel}
\usepackage[T1]{fontenc}
\usepackage[utf8]{inputenc}
\usepackage{fancyhdr}
\usepackage{float}
\usepackage{graphicx}
\usepackage{wrapfig}
\usepackage{siunitx} %per scrivere il simbolo °
\usepackage{verbatim} %per i commenti1
\usepackage{subfig}
\usepackage{amsmath}
\usepackage{algorithm}
\usepackage{algpseudocode}
\setcounter{secnumdepth}{3}
\setcounter{tocdepth}{6}
\usepackage{multirow}
\newcommand{\minitab}[2][l]{\begin{tabular}#1 #2\end{tabular}}
\usepackage{rotating}
\usepackage{xfrac}
\usepackage{enumitem}
\usepackage{amsmath}

\DeclareMathOperator*{\argmax}{arg\,max}
\DeclareMathOperator*{\argmin}{arg\,min}

%\usepackage{booktabs,array}
%\usepackage{tikz}

%\usepackage{tabularx}

%\usepackage{chngcntr}
%\counterwithin{table}{section}

%------------------------------ colors
\usepackage[usenames,dvipsnames,table]{xcolor} % use colors on table and more
\definecolor{333}{RGB}{51, 51, 51} % define custom color
\definecolor{background}{RGB}{255, 254, 213}
\definecolor{comment}{RGB}{17,167,5}
\definecolor{keyword}{RGB}{195,47,8}
\definecolor{string}{RGB}{142,195,0}
\definecolor{number}{RGB}{90,84,84}
\definecolor{identifier}{RGB}{0,90,201}

%------------------------------ source code
\usepackage{listings}

\lstset{
  basicstyle=\footnotesize\sffamily,
  commentstyle=\itshape\color{gray},
  captionpos=b,
  frame=shadowbox,
  language=HTML,
  rulesepcolor=\color{333},
  tabsize=2
}

\lstdefinestyle{code}{
  backgroundcolor=\color{background},
  basicstyle=\footnotesize\sffamily,
  commentstyle=\color{comment},
  frame=L,
  identifierstyle=\color{identifier},
  keywordstyle=\color{keyword},
  numbers=left,
  numbersep=10pt,
  numberstyle=\tiny\color{number},
  stringstyle=\color{string},
  showstringspaces=false,  
  stepnumber=1,
  tabsize=2
}


%------------------------------ define Abstract environment, missing in the 'book' class
\newenvironment{abstract}{\cleardoublepage \null \vfill \begin{center}\bfseries\abstractname \end{center}}{\vfill\null}
\addto\captionsenglish{\renewcommand*\abstractname{Abstract}} % change Abstract title

%------------------------------ active url
\usepackage{url}
\renewcommand{\UrlFont}{\color{black}\small\ttfamily}
\usepackage[colorlinks=true, linkcolor=black, citecolor=black, urlcolor=black]{hyperref} % active ref
%------------------------------ macros
\newcommand{\sectionname}{Section} % define Section ref
\newcommand{\subsectionname}{Sub-section} % define Sub-section ref
\renewcommand*\arraystretch{1.4} % tables padding

%acronimi
\usepackage[printonlyused]{acronym}

\begin{document}
\frontmatter

\begin{titlepage} %------------------------------ TITLE PAGE
\begin{center}

\hspace{0.5cm}
\begin{minipage}{.20\textwidth}
  \includegraphics[height=2.5cm]{Images/UNIPD}
\end{minipage}\begin{minipage}{.90\textwidth}
  \begin{table}[H]
  \begin{tabular}{l}
  \scshape{\Large{\bfseries{Università degli Studi di Padova}}} \\
  \hline \\
  \scshape{\Large{Facoltà di Ingegneria}} \\
  \end{tabular}
  \end{table}
\end{minipage}

\vspace{1cm}
\emph{\Large{Corso~di~Laurea~Magistrale~in~Ingegneria~Informatica}} \\
\vspace{3cm}
\scshape{\Large{Tesina di Ricerca Operativa 2}} \\
\end{center}

\vspace{1cm}
\begin{center}
\scshape{\Large{\bfseries{Travelling Salesman}}}\\
\scshape{\Large{\bfseries{Problem}}}
\end{center}

\vspace{3.5cm}

\begin{center}
\emph{Autori}
\vspace{0.2cm}
\begin{table}[h]
\centering
\begin{tabular}{rl}
\vspace{0.2cm}
{Raffaele Di Nardo Di Maio} & {1204879}\\
{Cristina Fabris} & {1205722}\\
\end{tabular}
\end{table}

\end{center}

\vfill
\begin{center}
\hspace{-0.2cm}
\line(1, 0){360}\\
\textsc{Anno Accademico 2019-2020}
\end{center}
\end{titlepage}

%\begin{comment}
\begingroup %------------------------------ CONTENTS
  \makeatletter
  \let\ps@plain\ps@empty
  \makeatother
  \tableofcontents  
  \clearpage
\endgroup
%\end{comment}

\makeatletter
\def\@makechapterhead#1{%
  \vspace*{50\p@}%
  {\parindent \z@ \raggedright \normalfont
    \vskip 20\p@
    \interlinepenalty\@M
    \Huge \bfseries #1\par\nobreak
    \vskip 40\p@
  }}
\makeatother


\mainmatter
\chapter{Introduzione}\label{intro}
La seguente trattazione analizza il Problema del Commesso Viaggiatore (Travelling Salesman Problem, TSP), che consiste nell'individuare un circuito hamiltoniano di costo minimo in un assegnato grafo orientato G=(V,A)\cite{TSP}. La formulazione matematica di tale problema è la seguente:
$$
x_{ij}=
\begin{cases}
1 & se\;l'arco\;(i,j)\in A\;viene\;scelto\;nel\;circuito\;ottimo\\
0 & altrimenti\\
\end{cases}
$$
\begin{align}
& min \underset{(i,j)\in A}\sum{c_{ij}\;x_{ij}} \\\notag \\
& \underset{(i,j)\in \delta^{-}(j)}\sum{\;x_{ij}} = 1 & \forall\;j\in V \\\notag \\
& \underset{(i,j)\in \delta^{+}(i)}\sum{\;x_{ij}} = 1 & \forall\;i\in V \\\notag \\
& \underset{(i,j)\in \delta^{+}(S)}\sum{\;x_{ij}} \geq 1 & S \subset V\; :\; 1\in S\\\notag \\
& x_{ij}\geq 0\;\; intero & (i,j)\in A.\\\notag
\end{align}
Tuttavia le soluzioni algoritmiche presentante risolvono una sua variante, detta simmetrica, che viene applicata ad un grafo completo non orientato G=(V,E).\\ Di seguito viene riportata la formulazione matematica di tale versione:\\
\begin{align}
& min \underset{e\in E}\sum{c_e\;x_e}\\
& \underset{e\in \delta(v)}\sum{\;x_e} = 2 & \forall\;v\in V \\
& \underset{e\in E(S)}\sum{\;x_e} \leq \vert S\vert - 1 & \forall\;S\underset{\neq} \subset V: \vert S\vert\geq 3.\\\notag
\end{align}
A livello commerciale esistono diverse tipologie di risolutori di problemi di programmazione lineare intera, basati sul Branch \& Bound. I più conosciuti in circolazione sono i seguenti:
\begin{itemize}
\item{\textbf{IBM ILOG CPLEX Optimization Studio}\cite{ILOG}\\
è un soluzione analitica, sviluppata dall'IBM e gratuita a livello accademico.}
\item{\textbf{FICO® Xpress Optimization}\cite{FICO}\\
è stato prodotto dalla Fair Isaac Corporation(FICO) ed è costituito da 4 componenti principali: FICO® Xpress Insight, FICO® Xpress Executor, FICO® Xpress Solver e FICO® Xpress Workbench. Questa soluzione è disponibile gratuitamente solo nella versione Community, in cui però vengono applicate restrizioni sul numero di righe e colonne del tableau, di token non lineari e di funzioni dell'utente.}
\item{\textbf{Gurobi}\cite{GUROBI}\\
è una soluzione, sviluppata dalla Gurobi Optimization, che viene rilasciata anche con una versione accademica.}
\item{\textbf{COIN Branch and Cut solver (CBC)}\cite{CBC}\\
è un risolutore MIP(mixed-integer program) open-source scritto in C++ e sviluppato dalla Computational Infrastructure for Operations Researc (COIN).}
\end{itemize}
Nel Capitolo \ref{CPLEX} vengono riportate diverse soluzioni mat-euristiche e non per il problema del Commesso Viaggiatore, che fanno uso di ILOG CPLEX.\\
In commercio, il più noto ed efficiente software per la risoluzione del TSP è Concorde, sviluppato in ANSI C e disponibile per l'uso in ambito accademico\cite{concorde}.\\
Nel Capitolo \ref{HEURISTIC} vengono analizzati gli algoritmi euristici, sviluppati senza far uso di ILOG CPLEX.\\
Nel Capitolo \ref{PERF_PROF} vengono invece riportati i confronti, a livello temporale e di costo, delle soluzioni ottenute con i differenti algoritmi enunciati.\\
Nell'Appendice \ref{TSPlib}, \ref{CPLEX_func}, \ref{gnuplot}, \ref{perf_profile_py} vengono descritte rispettivamente la tipologia di istanze utilizzate, la documentazione utilizzata ed il funzionamento di CPLEX, il programma GNUPLOT utilizzato nella stampa delle soluzioni e il programma perfprof.py usato per creare i performance profile del Capitolo \ref{PERF_PROF}.\\
Tutti i tempi di esecuzione e i costi delle soluzioni, ottenuti mediante la fase di testing, sono consultabili invece nelle tabelle riportate nell'Appendice \ref{results}. Tutte le soluzioni descritte sono state implementate in linguaggio C e testate sul sistema operativo Windows 10 con Visual Studio, ed i sorgenti sono disponibili online\footnote{\url{https://github.com/RaffaDNDM/Operational-Research-2}}.
\begin{thebibliography}{1}

\begin{comment}
\bibitem{alg} Pierluigi Crescenzi, Pierre Fraigniaud, Magnus Halldórsson, Hovhannes A. Harutyunyand, Chiara Pierucci, Andrea Pietracaprina, Geppino Pucci \emph {On the complexity of the shortest-path broadcast problem, Discrete Applied Mathematics , vol. 199, pp. 101-109, 2016.}
\end{comment}

\bibitem{cook} \emph{\url{http://www.math.uwaterloo.ca/tsp/}}

\bibitem{concorde} \emph{\url{http://www.math.uwaterloo.ca/tsp/concorde/index.html}}

\bibitem{monnalisa} \emph{\url{http://www.math.uwaterloo.ca/tsp/data/ml/monalisa.html}}

\bibitem{salvagnin_perf} Dolan, E. D., Mor{\'e}, J.J., \emph{Benchmarking optimization software with performance profiles}, Mathematical Programming, vol. 91, n. 2, pp. 201--213, 2002, Springer


\end{thebibliography}

\end{document}